\documentclass[11pt]{article}

\usepackage[a4paper,margin=30mm]{geometry}
\usepackage{amsmath,amssymb,amsthm,mathtools}

\newtheorem{lemma}{Lemma}
\newtheorem{theorem}[lemma]{Theorem}
\newtheorem{corollary}[lemma]{Corollary}
\newtheorem{definition}[lemma]{Definition}
\newtheorem{remark}[lemma]{Remark}

\newcommand{\OptCost}{C_{\alpha}^{*}(I)}
\newcommand{\WorstCost}{C_{\alpha}^{W}(I)}
\newcommand{\LamI}{\lambda_{\alpha}(I)}

\begin{document}


\begin{lemma}[Largest Job is single in opt]
Let $S^*$ be an optimal schedule, and suppose that a job
$j^*$ of processing time $p$ is a singleton. Then every job $k$
with
\[
    p_k>p,
\]
is also a singleton.

Consequently, if $j^*$ is the only singleton job in $S^*$,
then $j^*$ is a largest job.
\end{lemma}

\begin{proof}
Let $M_0$ be the singleton machine containing $j^*$, so that its
load is $p$.  Consider any job $k$ with processing time
\[
    p_k>p,
\]
and let $M_i$ be the machine containing $k$.  Denote the load of
$M_i$ by $L_i$.

Exchange $j^*$ and $k$.  Set
\[
    \delta:=p_k-p>0.
\]
Before the exchange, the two machine loads are $(p,L_i)$.  After the
exchange, they are $(p_k,L_i-\delta)$.
All other machine loads remain unchanged.  Hence the change in cost is
\begin{align*}
    p_k^2+(L_i-\delta)^2-p^2-L_i^2
      &=2\delta(p_k-L_i).
\end{align*}
Because $L_i\geq p_k$, this quantity is nonpositive.  If $M_i$
contains any job besides $k$, then all processing times are positive
and the exchange would strictly decrease the cost of the schedule which 
contradicts its optimality. 
Note that the exchange might not directly result in a list schedule. 
But, rescheduling the jobs that finish last in this setup will only further 
reduce machine loads, so this does not contradict the cost reduction.
Or expressed differntly, at least one of the longest jobs must be single-scheduled.
Finally, every job larger than $p$ must also be single scheduled in $S^*$.
\end{proof}

\section{The perturbation lemma}

\begin{definition}[Perturbed Instance]
Let $I$ be an instance of
$P\parallel\sum_j p_jC_j(\alpha)$, and let $j^*$ be a job with
processing time $p_{j^*}=p $. For \(\varepsilon\geq 0\), define
\(I(\varepsilon)\) to be the instance obtained from \(I\) by increasing
only the processing time of \(j^*\) by \(\varepsilon\). Formally,
\[
p_j(\varepsilon)
:=
\begin{cases}
p+\varepsilon, & \text{if } j=j^*,\\
p_j,           & \text{if } j\neq j^*.
\end{cases}
\]
We call $I(\varepsilon)$ the perturbed instance of $I$.
\end{definition}

\begin{lemma}[Perturbing single job]
\label{lem:critical-singleton-perturbation}
Let $\alpha\in(0,\tfrac12]$, and let $j^*$ be a job of processing
time $p$.  Suppose that
\begin{enumerate}
    \item $j^*$ is single scheduled in an optimal schedule of I
    \item $j^*$ is not single scheduled in a obective maximizing schedule $S^W$ of $I$
    \item $j^*$ is last on its machine in $S^W$.
\end{enumerate}
Let $M_{W}(j^*)$ be the \(\mathrm{WSPT}\) machine containing
$j^*$. As $j^*$ is not single scheduled in $\sigma^{W}$, 
Thus, $S^W_{j^*}>0$, and the load of the WLPT machine containing
$j^*$ is $p+S^W_{j^*}$.

Then, for every $\varepsilon\geq0$,
\begin{equation}
    \lambda_{\alpha}(I(\varepsilon))
    \geq
    \frac{
        \WorstCost+(S^W_{j^{*}}+2\alpha p)\varepsilon
        +\alpha\varepsilon^2
    }{
        \OptCost+2\alpha p\varepsilon
        +\alpha\varepsilon^2
    }.
    \label{eq:ratio-perturbation-bound}
\end{equation}
\end{lemma}

\begin{proof}
For any fixed schedule $S$ in which  ajob $j$ is last scheduled on
a machine, increasing its processing time $p$ to
$p+\varepsilon$ changes the assignment's cost by
\begin{align}
    \Delta(\sigma(I),\varepsilon)
      &=
      \frac12\bigl((p+\varepsilon+S_{j})^2-(p+S_{j})^2\bigr)
      \notag\\
      &\quad+
      \left(\alpha-\frac12\right)
      \bigl((p+\varepsilon)^2-p^2\bigr)
      \notag\\
      &=(S_{j}+2\alpha p)\varepsilon+\alpha\varepsilon^2.
    \label{eq:fixed-assignment-cost-change}
\end{align}
where $S$ denotes the starting time of $j$ which is of course equal to the load preceding $j$. 


Keep the optimal assignment from $I$ in which $j^*$ is 
single scheduled. This assignment is feasible for $I(\varepsilon)$.
As $j^*$ is single scheduled in an optimal schedule, $\sigma^{*}_{j^*}=0$. From (Queyranne reference here), shows that
its cost on $I(\varepsilon)$ is
\[
    \OptCost+2\alpha p\varepsilon+\alpha\varepsilon^2.
\]
Obviously, 
\begin{equation}
    C_{\alpha}^{*}(I(\varepsilon))
    \leq
    \OptCost+2\alpha p\varepsilon+\alpha\varepsilon^2.
    \label{eq:opt-candidate-bound}
\end{equation}

Next, fix a list producing a schedule $S^W$ that has properties from the statement.
Increasing the processing time of $j^*$ leaves all assignment
decisions before $j^*$ unchanged.  After $j^*$ is assigned,
only the availability time of its machine increases.  Since
$j^*$ is last on that machine, every later job in the fixed list
was assigned to another machine; making the machine of $j^*$
available later cannot change any of those choices.  Thus the same
list produces the same assignment for $I(\varepsilon)$.

Taking $S=S_{j^*}$ in \eqref{eq:fixed-assignment-cost-change} gives the
cost
\[
    \WorstCost+(S_{j^*}+2\alpha p)\varepsilon
    +\alpha\varepsilon^2.
\]
It remains a list schedule for $I(\varepsilon)$, whereas
$C_{\alpha}^{W}(I(\varepsilon))$ is the largest cost of any list schedule (for $I(\varepsilon)$)
Consequently,
\begin{equation}
    C_{\alpha}^{W}(I(\varepsilon))
    \geq
    \WorstCost+(T_{j^*}+2\alpha p)\varepsilon
    +\alpha\varepsilon^2.
    \label{eq:wlpt-candidate-bound}
\end{equation}

Dividing \eqref{eq:wlpt-candidate-bound} by
\eqref{eq:opt-candidate-bound} gives
\eqref{eq:ratio-perturbation-bound}.
\end{proof}

\section{Improving the approximation ratio}

\begin{theorem}[Strict improvement when the accompanying load is large]
\label{thm:strict-perturbation-improvement}
Let $\alpha\in(0,\tfrac12]$, and suppose that
\begin{equation}
    \LamI\leq 1+\frac{1}{6\alpha}
    \label{eq:rho-upper-bound}
\end{equation}
for every instance $I$ of
$P\parallel\sum_jp_jC_j(\alpha)$.  Let $I$ and $j^*$ satisfy the
assumptions of
\ref{lem:critical-singleton-perturbation}.  If
\[
    \sigma^{W}_{j^*}>\frac p3,
\]
then
\begin{equation}
    \boxed{\lambda_{\alpha}(I(\varepsilon))>\LamI}
    \label{eq:strict-ratio-improvement}
\end{equation}
for every
\begin{equation}
    0<\varepsilon<
    6\left(\sigma^{W}_{j^*}-\frac p3\right).
    \label{eq:epsilon-improvement-range}
\end{equation}
\end{theorem}

\begin{proof}
By
\ref{lem:critical-singleton-perturbation}, it is enough to show
\[
    \frac{
        \WorstCost+(\sigma^{W}_{j^*}+2\alpha p)\varepsilon
        +\alpha\varepsilon^2
    }{
        \OptCost+2\alpha p\varepsilon
        +\alpha\varepsilon^2
    }
    >\LamI.
\]
Using $\WorstCost=\LamI\OptCost$ and rearanging terms, shows
that this strict inequality is equivalent to
\begin{equation}
    \sigma^{W}_{j^*}>
    (\LamI-1)(2\alpha p+\alpha\varepsilon).
    \label{eq:exact-improvement-condition}
\end{equation}
\eqref{eq:rho-upper-bound} gives
\[
    (\lambda-1)(2\alpha p+\alpha\varepsilon)
    \leq
    \frac{1}{6\alpha}
    (2\alpha p+\alpha\varepsilon)
    =
    \frac p3+\frac{\varepsilon}{6}.
\]
If \eqref{eq:epsilon-improvement-range} holds, then
\[
    \frac p3+\frac{\varepsilon}{6}<\sigma^{W}_{j^*}.
\]
Thus \eqref{eq:exact-improvement-condition} holds, and
\eqref{eq:ratio-perturbation-bound} yields
\eqref{eq:strict-ratio-improvement}.
\end{proof}
From \ref{thm:strict-perturbation-improvement} we can follow, that a $\lambda$-maximzing instance must fulfill
\[
S_{j^*}^{W} \le \frac{p}{3}
\]
\end{document}
